% Options for packages loaded elsewhere
\PassOptionsToPackage{unicode}{hyperref}
\PassOptionsToPackage{hyphens}{url}
\documentclass[
  ignorenonframetext,
]{beamer}
\newif\ifbibliography
\usepackage{pgfpages}
\setbeamertemplate{caption}[numbered]
\setbeamertemplate{caption label separator}{: }
\setbeamercolor{caption name}{fg=normal text.fg}
\beamertemplatenavigationsymbolsempty
% remove section numbering
\setbeamertemplate{part page}{
  \centering
  \begin{beamercolorbox}[sep=16pt,center]{part title}
    \usebeamerfont{part title}\insertpart\par
  \end{beamercolorbox}
}
\setbeamertemplate{section page}{
  \centering
  \begin{beamercolorbox}[sep=12pt,center]{section title}
    \usebeamerfont{section title}\insertsection\par
  \end{beamercolorbox}
}
\setbeamertemplate{subsection page}{
  \centering
  \begin{beamercolorbox}[sep=8pt,center]{subsection title}
    \usebeamerfont{subsection title}\insertsubsection\par
  \end{beamercolorbox}
}
% Prevent slide breaks in the middle of a paragraph
\widowpenalties 1 10000
\raggedbottom
\AtBeginPart{
  \frame{\partpage}
}
\AtBeginSection{
  \ifbibliography
  \else
    \frame{\sectionpage}
  \fi
}
\AtBeginSubsection{
  \frame{\subsectionpage}
}
\usepackage{iftex}
\ifPDFTeX
  \usepackage[T1]{fontenc}
  \usepackage[utf8]{inputenc}
  \usepackage{textcomp} % provide euro and other symbols
\else % if luatex or xetex
  \usepackage{unicode-math} % this also loads fontspec
  \defaultfontfeatures{Scale=MatchLowercase}
  \defaultfontfeatures[\rmfamily]{Ligatures=TeX,Scale=1}
\fi
\usepackage{lmodern}
\ifPDFTeX\else
  % xetex/luatex font selection
\fi
% Use upquote if available, for straight quotes in verbatim environments
\IfFileExists{upquote.sty}{\usepackage{upquote}}{}
\IfFileExists{microtype.sty}{% use microtype if available
  \usepackage[]{microtype}
  \UseMicrotypeSet[protrusion]{basicmath} % disable protrusion for tt fonts
}{}
\makeatletter
\@ifundefined{KOMAClassName}{% if non-KOMA class
  \IfFileExists{parskip.sty}{%
    \usepackage{parskip}
  }{% else
    \setlength{\parindent}{0pt}
    \setlength{\parskip}{6pt plus 2pt minus 1pt}}
}{% if KOMA class
  \KOMAoptions{parskip=half}}
\makeatother
\usepackage{color}
\usepackage{fancyvrb}
\newcommand{\VerbBar}{|}
\newcommand{\VERB}{\Verb[commandchars=\\\{\}]}
\DefineVerbatimEnvironment{Highlighting}{Verbatim}{commandchars=\\\{\}}
% Add ',fontsize=\small' for more characters per line
\newenvironment{Shaded}{}{}
\newcommand{\AlertTok}[1]{\textcolor[rgb]{1.00,0.00,0.00}{\textbf{#1}}}
\newcommand{\AnnotationTok}[1]{\textcolor[rgb]{0.38,0.63,0.69}{\textbf{\textit{#1}}}}
\newcommand{\AttributeTok}[1]{\textcolor[rgb]{0.49,0.56,0.16}{#1}}
\newcommand{\BaseNTok}[1]{\textcolor[rgb]{0.25,0.63,0.44}{#1}}
\newcommand{\BuiltInTok}[1]{\textcolor[rgb]{0.00,0.50,0.00}{#1}}
\newcommand{\CharTok}[1]{\textcolor[rgb]{0.25,0.44,0.63}{#1}}
\newcommand{\CommentTok}[1]{\textcolor[rgb]{0.38,0.63,0.69}{\textit{#1}}}
\newcommand{\CommentVarTok}[1]{\textcolor[rgb]{0.38,0.63,0.69}{\textbf{\textit{#1}}}}
\newcommand{\ConstantTok}[1]{\textcolor[rgb]{0.53,0.00,0.00}{#1}}
\newcommand{\ControlFlowTok}[1]{\textcolor[rgb]{0.00,0.44,0.13}{\textbf{#1}}}
\newcommand{\DataTypeTok}[1]{\textcolor[rgb]{0.56,0.13,0.00}{#1}}
\newcommand{\DecValTok}[1]{\textcolor[rgb]{0.25,0.63,0.44}{#1}}
\newcommand{\DocumentationTok}[1]{\textcolor[rgb]{0.73,0.13,0.13}{\textit{#1}}}
\newcommand{\ErrorTok}[1]{\textcolor[rgb]{1.00,0.00,0.00}{\textbf{#1}}}
\newcommand{\ExtensionTok}[1]{#1}
\newcommand{\FloatTok}[1]{\textcolor[rgb]{0.25,0.63,0.44}{#1}}
\newcommand{\FunctionTok}[1]{\textcolor[rgb]{0.02,0.16,0.49}{#1}}
\newcommand{\ImportTok}[1]{\textcolor[rgb]{0.00,0.50,0.00}{\textbf{#1}}}
\newcommand{\InformationTok}[1]{\textcolor[rgb]{0.38,0.63,0.69}{\textbf{\textit{#1}}}}
\newcommand{\KeywordTok}[1]{\textcolor[rgb]{0.00,0.44,0.13}{\textbf{#1}}}
\newcommand{\NormalTok}[1]{#1}
\newcommand{\OperatorTok}[1]{\textcolor[rgb]{0.40,0.40,0.40}{#1}}
\newcommand{\OtherTok}[1]{\textcolor[rgb]{0.00,0.44,0.13}{#1}}
\newcommand{\PreprocessorTok}[1]{\textcolor[rgb]{0.74,0.48,0.00}{#1}}
\newcommand{\RegionMarkerTok}[1]{#1}
\newcommand{\SpecialCharTok}[1]{\textcolor[rgb]{0.25,0.44,0.63}{#1}}
\newcommand{\SpecialStringTok}[1]{\textcolor[rgb]{0.73,0.40,0.53}{#1}}
\newcommand{\StringTok}[1]{\textcolor[rgb]{0.25,0.44,0.63}{#1}}
\newcommand{\VariableTok}[1]{\textcolor[rgb]{0.10,0.09,0.49}{#1}}
\newcommand{\VerbatimStringTok}[1]{\textcolor[rgb]{0.25,0.44,0.63}{#1}}
\newcommand{\WarningTok}[1]{\textcolor[rgb]{0.38,0.63,0.69}{\textbf{\textit{#1}}}}
\usepackage{longtable,booktabs,array}
\usepackage{caption}
\captionsetup[table]{skip=6pt}
\newcounter{none} % for unnumbered tables
\usepackage{calc} % for calculating minipage widths
\usepackage{caption}
% Make caption package work with longtable
\makeatletter
\def\fnum@table{\tablename~\thetable}
\makeatother
\setlength{\emergencystretch}{3em} % prevent overfull lines
\providecommand{\tightlist}{%
  \setlength{\itemsep}{0pt}\setlength{\parskip}{0pt}}
% Auto-generated by infrastructure/rendering/_pdf_latex_helpers.py::extract_math_font_preamble.
% Minimal header passed to Pandoc via -H so Beamer slides pick up the same math font
% as the combined-PDF path. Adjust the manuscript's preamble.md to change the font globally.
\usepackage{unicode-math}
\setmathfont{latinmodern-math.otf}

% Slide layout defaults for warning-clean scientific decks.
\usepackage{etoolbox}
\IfFileExists{xurl.sty}{\usepackage{xurl}}{}
\IfFileExists{seqsplit.sty}{\usepackage{seqsplit}}{\newcommand{\seqsplit}[1]{#1}}
\protected\def\breakseq#1{\seqsplit{#1}}
\protected\def\breaktt#1{\begingroup\ttfamily\seqsplit{#1}\endgroup}
\setlength{\emergencystretch}{6em}
\tolerance=5000
\hbadness=10000
\hfuzz=1pt
\setlength{\tabcolsep}{2pt}
\AtBeginEnvironment{longtable}{\tiny\renewcommand{\arraystretch}{0.86}\setlength{\tabcolsep}{1pt}}
\AtBeginEnvironment{tabular}{\tiny\renewcommand{\arraystretch}{0.86}\setlength{\tabcolsep}{1pt}}
\AtBeginEnvironment{equation}{\tiny}
\AtBeginEnvironment{equation*}{\tiny}
\AtBeginEnvironment{align}{\tiny}
\AtBeginEnvironment{align*}{\tiny}
\AtBeginEnvironment{itemize}{\footnotesize}
\AtBeginEnvironment{enumerate}{\footnotesize}
\AtBeginEnvironment{description}{\footnotesize}
\setbeamerfont{caption}{size=\tiny}
\setbeamerfont{caption name}{size=\tiny}
\setbeamerfont{normal text}{size=\small}
\setbeamerfont{frametitle}{size=\small}
\setbeamerfont{section title}{size=\footnotesize}
\setbeamerfont{subsection title}{size=\footnotesize}
\setbeamertemplate{section page}{%
  \centering
  \begin{beamercolorbox}[sep=12pt,center,wd=\paperwidth]{section title}
    \parbox{0.86\paperwidth}{\centering\usebeamerfont{section title}\insertsection\par}
  \end{beamercolorbox}
}
\setbeamertemplate{subsection page}{%
  \centering
  \begin{beamercolorbox}[sep=8pt,center,wd=\paperwidth]{subsection title}
    \parbox{0.86\paperwidth}{\centering\usebeamerfont{subsection title}\insertsubsection\par}
  \end{beamercolorbox}
}
\setlength{\abovecaptionskip}{2pt}
\setlength{\belowcaptionskip}{0pt}

% Natbib and cross-reference fallbacks — slides don't load natbib
% or cleveref, but combined-PDF manuscript prose may emit these
% commands. The fallback renders citations as a bracketed key list
% and cross-references as detokenized labels so slides don't fail on
% undefined control sequences or raw underscores. \providecommand is
% a no-op if packages load later.
\providecommand{\citep}[1]{[#1]}
\providecommand{\citet}[1]{#1}
\providecommand{\citealp}[1]{#1}
\providecommand{\citeauthor}[1]{#1}
\providecommand{\citeyear}[1]{#1}
\providecommand{\cref}[1]{\texttt{\detokenize{#1}}}
\providecommand{\Cref}[1]{\texttt{\detokenize{#1}}}

\newtheorem{proposition}{Proposition}
\newtheorem{hypothesis}{Hypothesis}
\usepackage{bookmark}
\IfFileExists{xurl.sty}{\usepackage{xurl}}{} % add URL line breaks if available
\urlstyle{same}
% fallback for those not using the hyperref driver hyperxmp:
\makeatletter
\@ifundefined{xmpquote}{\newcommand{\xmpquote}[1]{#1}}{}
\makeatother
\hypersetup{
  hidelinks,
  pdfcreator={LaTeX via pandoc}}

\author{\texorpdfstring{}{}}
\date{}

\begin{document}

\section{Evaluation Setup: Deterministic Artifact Tests and Planned
Learner
Study}\label{sec:evaluation-setup-deterministic-tests-learner-study}

\subsection{Deterministic model suite and fixed
inputs}\label{deterministic-model-suite-and-fixed-inputs}

\begin{frame}[fragile,allowframebreaks]{Deterministic model suite and
fixed inputs}
All quantitative results derive from small, fully specified discrete
models evaluated by the kernels with a fixed seed (\texttt{20260609}).
Within the supported Python, NumPy and Matplotlib environment, the model
inputs and graph ordering are deterministic; this is a
software-environment reproducibility claim, not a cross-hardware bitwise
proof.

\end{frame}

\begin{frame}[fragile,allowframebreaks]{Deterministic model suite and
fixed inputs}
\begin{itemize}
\tightlist
\item
  \textbf{Noisy-sensor model} (Bayes and VFE results): two latent
  states, \texttt{present} and \texttt{absent}, with prior
  \(D=[0.01, 0.99]\) and likelihood rows \([0.99, 0.05]\) and
  \([0.01, 0.95]\) over \texttt{positive} and \texttt{negative}
  observations. These model parameters are registered as kernel claims,
  so the prose cannot drift from \breaktt{src/kernels/examples.py}. This
  model emits the posterior, evidence probability, negative
  log-evidence, and free-energy sweep.
\item
  \textbf{Informative-sensor model} (EFE decomposition): likelihood rows
  \([0.9, 0.1]\) and \([0.1, 0.9]\) with observation preference
  \(C=[0.9, 0.1]\). This isolates epistemic value under uncertain versus
  certain beliefs.
\item
  \textbf{Two-action agent}: a neutral-preference sample/commit model
  with prior \(D=[0.5,0.5]\). The \emph{sample} action preserves the
  uncertain belief; the \emph{commit} action collapses belief mass
  toward one state. A sensor-precision sweep over this same model
  identifies the sample/commit switch.
\item
  \textbf{Information examples}: a fair coin, a fair six-sided die, a
  biased \(0.9/0.1\) coin, a one-in-a-hundred event, and a correlated
  joint distribution with entries \([0.4,0.1,0.1,0.4]\).
\end{itemize}
\end{frame}

\subsection{Curriculum export and runtime ownership
boundary}\label{curriculum-export-and-runtime-ownership-boundary}

\begin{frame}[fragile,allowframebreaks]{Curriculum export and runtime
ownership boundary}
The curriculum spans the eight strata declared in
\breaktt{manuscript/config.yaml}. The maximal-enumeration floor is 630
skills (achieved: 630), and the current catalog adds
practice-and-translation material across all strata. That material
specifies opportunities to rehearse model checking, debugging, policy
design, implementation review, and application evaluation; it does not
establish that learners will acquire those abilities. Skills default to
100 points, one occurrence to completion, and \texttt{Quiz}
self-reporting so completion is assessment-gated. SkillTree export
targets project id \texttt{activeInference}.
\end{frame}

\subsection{Reproduction, delivery, and release
checks}\label{reproduction-delivery-and-release-checks}

\begin{frame}[fragile,allowframebreaks]{Reproduction, delivery, and
release checks}
The generation pipeline is deterministic and runs from the project root
with no Docker daemon and no browser. Only
\breaktt{validate\_course\_material\_sources.py\ -\/-online} leaves the
machine, and only to re-resolve the cited course-material sources:

\begin{Shaded}
\begin{Highlighting}[]
\ExtensionTok{uv}\NormalTok{ run python scripts/build\_curriculum.py   }\CommentTok{\# validate the DAG, emit stats}
\ExtensionTok{uv}\NormalTok{ run python scripts/run\_kernels.py        }\CommentTok{\# claims, scholarship, figures, artifact audit}
\ExtensionTok{uv}\NormalTok{ run python scripts/build\_course\_materials.py}
\ExtensionTok{uv}\NormalTok{ run python scripts/validate\_course\_material\_sources.py }\AttributeTok{{-}{-}online}
\ExtensionTok{uv}\NormalTok{ run python scripts/validate\_manuscript\_artifacts.py}
\ExtensionTok{uv}\NormalTok{ run python scripts/validate\_docker\_digests.py}
\ExtensionTok{uv}\NormalTok{ run python scripts/validate\_release\_readiness\_inventory.py}
\ExtensionTok{uv}\NormalTok{ run python scripts/export\_skilltree.py   }\CommentTok{\# provenance{-}gated SkillTree JSON}
\ExtensionTok{uv}\NormalTok{ run pytest }\AttributeTok{{-}{-}cov}\OperatorTok{=}\NormalTok{src }\AttributeTok{{-}{-}cov{-}fail{-}under}\OperatorTok{=}\NormalTok{90 }\CommentTok{\# coverage gate}
\end{Highlighting}
\end{Shaded}

\end{frame}

\begin{frame}[fragile,allowframebreaks]{Reproduction, delivery, and
release checks}
\breaktt{validate\_docker\_digests.py} belongs in that offline list
because it reads \breaktt{docker-compose.yml} as a file rather than
talking to a daemon: it fails when any Compose image is named by a
mutable tag instead of a \texttt{sha256} digest, which is what keeps the
live rehearsal below reproducible instead of tag-dependent. Its
\texttt{-\/-file} flag points the same check at another Compose file.

The remaining checks are live. All of them need a reachable Docker
daemon, and the browser-facing ones additionally need the
\breaktt{chrome-devtools-axi} Chrome driver. They fail closed rather than
downgrading to the artifact-only dashboard:

\end{frame}

\begin{frame}[fragile,allowframebreaks]{Reproduction, delivery, and
release checks}
\begin{Shaded}
\begin{Highlighting}[]
\ExtensionTok{./scripts/run\_skilltree.sh}                  \CommentTok{\# docker compose up + seed}
\ExtensionTok{./scripts/run\_learning\_dashboard.sh}         \CommentTok{\# SkillTree + seed + local dashboard}
\ExtensionTok{uv}\NormalTok{ run python scripts/reconcile\_live\_dependencies.py }\AttributeTok{{-}{-}check{-}only}
\ExtensionTok{uv}\NormalTok{ run python scripts/validate\_live\_skilltree\_dashboard.py }\AttributeTok{{-}{-}require{-}live{-}admin} \AttributeTok{{-}{-}require{-}live{-}dependencies} \AttributeTok{{-}{-}require{-}figure{-}contracts}
\ExtensionTok{uv}\NormalTok{ run python scripts/validate\_skilltree\_form\_login.py}
\ExtensionTok{uv}\NormalTok{ run python scripts/validate\_dashboard\_console.py }\AttributeTok{{-}{-}require{-}browser} \AttributeTok{{-}{-}require{-}dashboard}
\ExtensionTok{uv}\NormalTok{ run python scripts/validate\_full\_app\_stack.py }\AttributeTok{{-}{-}app{-}ready} \AttributeTok{{-}{-}skip{-}regenerate}
\ExtensionTok{uv}\NormalTok{ run python scripts/validate\_full\_app\_stack.py }\AttributeTok{{-}{-}skip{-}regenerate}
\end{Highlighting}
\end{Shaded}

\end{frame}

\begin{frame}[fragile,allowframebreaks]{Reproduction, delivery, and
release checks}
\breaktt{run\_learning\_dashboard.sh} regenerates the curriculum, claims,
course-material index, figure registry, manuscript artifact audit,
figures, and SkillTree export before seeding, so the learner surface is
inspected as an actual course interface rather than a static preview.
With local auth enabled, the dashboard serves
\breaktt{http://127.0.0.1:8765/api/skilltree/token} as a trusted-client
bridge for the official SkillTree display; it never embeds the project
client secret in static JavaScript. The dashboard can also run in
artifact-only mode with
\breaktt{uv\ run\ python\ scripts/serve\_learner\_dashboard.py\ -\/-check},
which verifies that the server, route map, and artifact inventory load
without requiring Docker.

\end{frame}

\begin{frame}[fragile,allowframebreaks]{Reproduction, delivery, and
release checks}
The live gates divide by what they are willing to assert.
\breaktt{reconcile\_live\_dependencies.py\ -\/-check-only} reports live
prerequisite parity against \breaktt{output/data/skilltree\_project.json}
without writing anything; dropping \texttt{-\/-check-only} writes the
missing edges into the running service, with
\breaktt{-\/-attempts-per-edge}, \texttt{-\/-max-edges}, and
\breaktt{-\/-max-runtime-seconds} bounding a chunked local run.
\breaktt{validate\_live\_skilltree\_dashboard.py} checks the SkillTree
and dashboard status endpoints, export-count agreement, artifact
completeness, and the trusted-client token exchange by default; its
\breaktt{-\/-require-live-admin}, \breaktt{-\/-require-live-dependencies},
and \breaktt{-\/-require-figure-contracts} flags escalate that to live
admin skill-count parity, exact live dependency parity, and
figure-registry-to-dashboard contract agreement.
\end{frame}

\begin{frame}[fragile,allowframebreaks]{Reproduction, delivery, and
release checks}
\breaktt{validate\_skilltree\_form\_login.py} drives the raw SkillTree
\texttt{FORM} login in a real browser through to the administrator
surface and redacts local credentials from its JSON summary.
\breaktt{validate\_dashboard\_console.py} with
\breaktt{-\/-require-browser\ -\/-require-dashboard} fails on any
dashboard console error other than the allow-listed upstream
\texttt{Inception} status probe; without those flags it reports a skip
when the browser or dashboard is absent.

\end{frame}

\begin{frame}[fragile,allowframebreaks]{Reproduction, delivery, and
release checks}
\breaktt{validate\_full\_app\_stack.py} composes the live sequence end to
end: Compose bring-up, seeding, dashboard start, dependency
reconciliation, the live dashboard checks, the form login, and the
console gate. \texttt{-\/-app-ready} skips dependency reconciliation and
labels its JSON summary app-ready; the flagless run additionally
requires every exported prerequisite edge to be present in the live
service and labels the summary full-parity. \breaktt{-\/-skip-regenerate}
reuses existing \texttt{output/} artifacts instead of rerunning the
generation sequence, \breaktt{-\/-skip-browser} drops the browser stage,
and \breaktt{-\/-keep-running} leaves the dashboard up for manual
inspection. Each successful run writes a redacted JSON summary under the
local state directory. That evidence is per-run and host-local: it
records what one rehearsal observed on one machine, not a public
deployment.
\end{frame}

\subsection{Planned learner-outcome study (design
only)}\label{planned-learner-outcome-study-design-only}

\begin{frame}[fragile,allowframebreaks]{Planned learner-outcome study
(design only)}
The next learner-outcome step is preregistered in
\breaktt{experiment\_plan.yaml} as design-only work, not as a result. The
planned pilot uses adult volunteer consent, data minimization, and
deidentified study ids before any SkillTree-owned progress export is
analyzed. Learner progress, quiz attempts, completion, and ranking
remain owned by SkillTree; the local dashboard is not treated as a
learner-truth database.

The planned measurement sequence is baseline concept items, an assigned
SkillTree learning path, immediate post-items, delayed-retention items,
and an optional transfer prompt scored by a frozen rubric.
Active-learning, cognitive-engagement, retrieval-practice, and
evidence-centered assessment sources motivate that design envelope, but
they do not substitute for data from this curriculum (Freeman et al.
2014; Chi and Wylie 2014; Roediger and Karpicke 2006; Mislevy et al.
2003; Black and Wiliam 1998; Shute 2008; Monib et al. 2025; Rof et al.
2024; Emin et al. 2026).
\end{frame}

\begin{frame}[fragile,allowframebreaks]{Planned learner-outcome study
(design only)}

Design-based research is the intended study posture because the
intervention is a curriculum artifact, a SkillTree delivery
configuration, and a measurement workflow that must be studied in
context and iterated before stronger outcome language would be
appropriate (Design-Based Research Collective 2003; Barab and Squire
2004). Learning-analytics reviews and intervention meta-analyses
motivate the outcome-evaluation path but also make privacy, data
protection, and instrumentation validity part of the study boundary (Liu
et al. 2025; Hernández-Campos et al. 2025; Liu and Khalil 2023; Sclater
and Bailey 2015). The protocol will use pilot-and-feasibility reporting
and an intervention-description checklist so that implementation
readiness is reported separately from definitive learner outcomes
(Eldridge et al. 2016; Hoffmann et al. 2014). The analysis plan is
within-learner baseline-to-post change with delayed retention reported
\end{frame}

\begin{frame}[fragile,allowframebreaks]{Planned learner-outcome study
(design only)}
separately after the sample size, item families, exclusions, and rubric
are frozen. This manuscript does not report those learner data and makes
no efficacy, retention, or transfer claim.
\end{frame}

\begin{frame}[fragile,allowframebreaks]{Browser-facing validation path}
\protect\phantomsection\label{browser-facing-validation-path}
For end-to-end browser validation, open the dashboard, confirm the
official SkillTree panel mounts the \texttt{activeInference} project,
switch to the curriculum panel, and open a skill's \textbf{Learn},
\textbf{Open Subject}, and \textbf{Quiz} actions. \textbf{Learn} opens
the local modular lesson record and should not emit a direct
\texttt{/subjects/\{subjectId\}/skills/\{skillId\}} SkillTree target.
\textbf{Open Subject} reopens the dashboard with
\texttt{skillsClientDisplayPath=/subjects/\{subjectId\}} so the subject
learning path, skill list, progress, and rank are exercised inside the
official SkillTree display. \textbf{Quiz} opens
\texttt{skillsClientDisplayPath=/subjects/\{subjectId\}/skills/\{skillId\}/quizzes/\{quizId\}},
so quiz loading, scoring rules, and completion remain SkillTree-owned
rather than dashboard-owned.
\end{frame}

\end{document}
